\chapter{Deep Dive: The Dictionary on MongoDB --- Document Modeling, Regex Search, and Honest Trade-offs}
\label{ch:dd-mongo}

Chapter~8 made the case for the dictionary as the project's document
service: an entry is read as a unit, written as a unit, and naturally
nested. That case is sound. This chapter tests it the way a skeptical
senior engineer would --- by reading the queries the service actually
runs, asking what each one costs at ten thousand words per user, and
admitting where a relational database would have done as well. The goal
is not to reverse the decision but to let you defend either side of it
in an interview, with the code in front of you.

Three questions organize the chapter:

\begin{enumerate}
  \item \textbf{What does the document model buy here}, and what would
  Postgres with \code{jsonb} have cost instead?
  \item \textbf{How does the search query actually execute} --- which
  part touches an index, which part scans, and where does the page get
  built?
  \item \textbf{What breaks when a dictionary gets big} --- and which of
  those breaks are Mongo's fault, and which are the code's?
\end{enumerate}

\section{Why MongoDB here, and what it solves}

Open \code{WordDefinition}. It is the whole schema of the service:

\begin{lstlisting}[language=Java]
@Document(collection = "word_definitions")
public class WordDefinition {
    @Id private String id;
    private String word;
    private String createdByUserId;                       (*@\co{1}@*)
    private String meaning;
    private String usage;
    private String notes;
    private List<String> examples;
    private List<String> tags;                            (*@\co{2}@*)
    /** Image attachments stored inline as base64 data URLs. */
    private List<String> images;                          (*@\co{3}@*)
    private Map<String, Object> extra;                    (*@\co{4}@*)
    @CreatedDate private Instant createdAt;
    @LastModifiedDate private Instant updatedAt;
}
\end{lstlisting}

\begin{description}
  \item[\co{1}] Every query in the service is scoped by this field. It
  is the tenant fence and, as we will see, the only thing the index
  helps with.
  \item[\co{2}] Two arrays embedded in the document. In Postgres these
  are two child tables or two \code{text[]} columns; here they are part
  of the entry and cost nothing to fetch with it.
  \item[\co{3}] Images as inline base64 strings. This is Chapter~35's
  ``bytes in the database'' in document clothing, with a hard ceiling
  attached: a BSON document may not exceed 16\,MB. A dozen phone photos
  and a word cannot be saved at all.
  \item[\co{4}] A free-form map. The document model's escape hatch:
  fields the code does not know about yet. Also the reason
  ``schemaless'' needs discipline (Chapter~8).
\end{description}

What the model buys: one read returns everything the page needs, no
joins, no N+1 across tables; adding a field is a code change, not a
migration; and the shape on the wire is the shape on disk, which makes
the service almost entirely glue. What it costs: relationships between
entries (\code{WordLink}) are not part of the document and must be
traversed by hand; every array field is un-indexable-by-prefix for
substring search; and nothing stops a document from growing without
bound.

\begin{decision}[MongoDB, and the case for having chosen Postgres]
The dictionary is a per-user notebook with nested, optional, evolving
fields --- the document model's home turf --- and it shares nothing
with the relational services, so isolation costs nothing. That is a
good reason. The honest counter: Postgres has \code{jsonb} for the
nested parts, \code{text[]} with a GIN index for tags,
\code{pg\_trgm} for exactly the substring search this service runs, and
recursive CTEs for the word graph, all in a database the team already
operates twice. A senior answer to ``why Mongo?'' is therefore not
``because documents''; it is ``because the data is document-shaped
\emph{and} we accepted a second database to keep the service
independently deployable and to learn polyglot persistence --- and here
is what it cost us''. What would flip it: a requirement to join words
to courses or users, or an operations team unwilling to run a second
engine.
\end{decision}

\section{Reading the search path}

The endpoint everyone uses is search. It runs through three pieces of
code. The repository query:

\begin{lstlisting}[language=Java]
@Query("{ 'createdByUserId': ?0, '$or': [ "
        + "{ 'word': { '$regex': ?1, '$options': 'i' } }, "
        + "{ 'tags': { '$regex': ?1, '$options': 'i' } } ] }")
List<WordDefinition> searchByWordOrTag(String userId, String regex);
\end{lstlisting}

The index that exists, from \code{MongoIndexConfig}:

\begin{lstlisting}[language=Java]
mongoTemplate.indexOps("word_definitions")
        .ensureIndex(new Index().on("createdByUserId", Sort.Direction.ASC));
\end{lstlisting}

And the service method that turns the list into a page:

\begin{lstlisting}[language=Java]
List<WordDefinitionResponse> responses = wordDefinitionRepository
        .searchByWordOrTag(userId, escapeRegex(query.trim()))   (*@\co{1}@*)
        .stream()
        .map(w -> toResponse(w, userId))                        (*@\co{2}@*)
        .sorted(Comparator.comparing(WordDefinitionResponse::word,
                String.CASE_INSENSITIVE_ORDER))                 (*@\co{3}@*)
        .toList();
int start = (int) pageable.getOffset();
int end = Math.min(start + pageable.getPageSize(), responses.size());
List<WordDefinitionResponse> pageContent = start < responses.size()
        ? responses.subList(start, end) : List.of();            (*@\co{4}@*)
return new PageImpl<>(pageContent, pageable, responses.size());
\end{lstlisting}

\begin{description}
  \item[\co{1}] The user's text is spliced into a regular expression.
  \code{escapeRegex} backslashes every metacharacter first, so
  \code{a.*} matches the three literal characters and \code{(} does not
  crash the server with a parse error. Skip this and you have regex
  injection: at best wrong results, at worst a catastrophic pattern
  that pins a CPU core.
  \item[\co{2}] \code{toResponse} runs two more queries per word ---
  parents and children from \code{word\_links} --- so a search that
  matches $n$ words issues $2n + 1$ round trips. This is the N+1
  problem, reborn in a database that was supposed to avoid it.
  \item[\co{3}] Sorting happens in the JVM, on \emph{every} match, not
  on the page.
  \item[\co{4}] So does paging. Mongo returned all matches; the service
  keeps twenty and discards the rest. Page 50 costs exactly what page 1
  costs, plus the copying.
\end{description}

\subsection{How the server executes it}

Ask Mongo itself. With the stack running:

\begin{handson}[read the plan]
\begin{lstlisting}[language=bash]
$ docker compose exec mongodb mongosh ts_dictionary --quiet --eval '
  db.word_definitions.find({ createdByUserId: "1", $or: [
    { word: { $regex: "clou", $options: "i" } },
    { tags: { $regex: "clou", $options: "i" } } ] })
  .explain("executionStats").executionStats'
\end{lstlisting}
Trimmed, the answer looks like this:
\begin{lstlisting}
{ nReturned: 3,
  totalKeysExamined: 412,
  totalDocsExamined: 412,
  executionStages: {
    stage: "FETCH",
    filter: { $or: [ { word: { $regex: "clou", $options: "i" } },
                     { tags: { $regex: "clou", $options: "i" } } ] },
    inputStage: {
      stage: "IXSCAN",
      indexName: "createdByUserId_1",
      indexBounds: { createdByUserId: [ "[\"1\", \"1\"]" ] } } } }
\end{lstlisting}
Three lines tell the story: \code{IXSCAN} on \code{createdByUserId}
narrows to this user's 412 documents; \code{FETCH} then loads every one
of them from disk; the \code{filter} runs both regexes against each. The
ratio \code{totalDocsExamined / nReturned} is 137 --- for each word
returned, 137 were read and rejected.
\end{handson}

Why the regex cannot use an index: a B-tree orders strings by their
leading characters. \code{\^{}clou} (anchored, case-sensitive) maps to a
key range and Mongo will use an index on \code{word} for it. \code{clou}
unanchored can appear anywhere in the string, and \code{i} means the
index's byte order does not even match the comparison; there is no range
to seek to. Mongo's planner is explicit about this: a case-insensitive
or unanchored \code{\$regex} is a full scan of whatever the other
predicates left. Here that is the user's slice, which is why the
\code{createdByUserId} index is doing all the useful work and the
compound index proposed later does not change the regex part --- it
changes what happens \emph{after} it.

\subsection{The array predicate}

\code{tags} is an array. A scalar condition on an array field matches if
\emph{any} element matches --- there is no \code{\$elemMatch} needed for
a single condition. That is a genuine convenience over SQL, where the
same predicate is an \code{EXISTS} against a child table. It is also
why the regex is run once per tag per document: a document with eight
tags costs eight pattern evaluations on top of the one for \code{word}.

\section{The graph}

\code{WordLink} is an adjacency list: one document per edge with
\code{parentWordId}, \code{childWordId}, \code{position}. The indexes
on both ends are right --- ``children of X'' and ``parents of X'' are
each one index lookup. The cost is in the traversal:

\begin{lstlisting}[language=Java]
private WordGraphResponse buildGraphNode(String wordId, String userId,
                                         Set<String> visited) {
    if (visited.contains(wordId)) return null;
    visited.add(wordId);
    WordDefinition word = wordDefinitionRepository
            .findByIdAndCreatedByUserId(wordId, userId).orElse(null); (*@\co{1}@*)
    if (word == null) return null;
    List<WordLink> childLinks = wordLinkRepository
            .findByParentWordIdOrderByPosition(wordId);              (*@\co{2}@*)
    List<WordGraphResponse> children = childLinks.stream()
            .map(link -> buildGraphNode(link.getChildWordId(), userId,
                    visited))
            .filter(Objects::nonNull).toList();
    return new WordGraphResponse(word.getId(), word.getWord(),
            word.getMeaning(), children);
}
\end{lstlisting}

\begin{description}
  \item[\co{1}] One query per node.
  \item[\co{2}] One more per node. A graph of 2{,}000 words costs
  4{,}000 round trips, sequentially, on the request thread. Each is
  sub-millisecond on the same host; on a network hop it is a
  multi-second page. Embedding \code{childIds} in the parent document
  would make it two queries total (\code{find} all words for the user,
  build in memory) --- at the price of updating two documents per link,
  without a transaction to hold them together.
\end{description}

Mongo does have a graph operator, \code{\$graphLookup}, that walks
\code{word\_links} recursively server-side in one aggregation. It
returns the flattened set of reachable nodes with a depth field, and
the tree is rebuilt in memory. The project's \code{visited} set is doing
by hand what \code{\$graphLookup} does with \code{maxDepth}; the
difference is $2n$ round trips versus one.

\section{Transactions and consistency, honestly}

Mongo supports multi-document ACID transactions since 4.0 --- on a
replica set. The project runs a single \code{mongod} (Compose and
StatefulSet alike), and a standalone node cannot start a transaction at
all. So \code{deleteWord} is two independent writes: delete the
definition, then \code{deleteByParentWordIdOrChildWordId}. A crash
between them leaves edges pointing to a missing word; the traversal
tolerates it (\code{orElse(null)}, then filtered), which is the right
defensive posture, but it is a hidden data-integrity contract. The
fix that costs nothing in code and one line in operations is to run a
single-node replica set (\code{--replSet rs0} plus an initiate call):
transactions become available, and so does the change stream you will
want for the outbox pattern of Chapter~11.

Write concern and read preference matter the moment there is a second
node. The default \code{w: majority} means an acknowledged write is on a
majority of nodes; reading from a secondary can return an older state.
For a single-user notebook, read-your-own-writes on the primary is what
users expect, and it is the default --- state that in an interview, and
that you would revisit it before adding read replicas.

\section{Pros and cons, candidly}

\begin{center}
\begin{tabular}{@{}p{0.3\linewidth}p{0.62\linewidth}@{}}
\toprule
\textbf{Strength} & \textbf{Where the project pays for it} \\
\midrule
One document, one read, no joins &
  Links are separate; the graph is $2n$ queries \\
Arrays are first-class (\code{tags}, \code{examples}) &
  Substring search over them is a scan per element \\
No migrations for new fields &
  \code{extra} and \code{images} grow unchecked toward 16\,MB \\
Regex search works out of the box &
  It cannot use an index unless anchored and case-sensitive \\
Paging is one \code{skip/limit} away &
  The search path pages in Java after loading everything \\
Independent database, independent deploy &
  No auth on \code{mongod} in Compose \emph{or} the cluster \\
\bottomrule
\end{tabular}
\end{center}

The last row is not a Mongo trait but a project fact worth knowing
before someone asks: the official image enables authentication only when
root credentials are set, and neither manifest sets them. Inside a
ClusterIP-only network that is a considered risk, documented in
\code{mongodb.yaml}; it is the first thing to change if the service ever
shares a cluster with anything untrusted.

\section{What breaks}

\begin{pitfall}[search gets slower with every word, then times out]
Symptom: search latency grows linearly with a user's dictionary; at a
few thousand words the frontend's timeout fires and the page shows
nothing. Cause: three multipliers stacked on the same request --- the
regex scans every document in the user's slice; \code{toResponse} then
issues two queries per match; and sorting and paging happen after both.
The database is doing what it was asked. Fix in order of payoff: page in
the database (below), batch the link lookups into one \code{\$in} query
per page, and anchor or replace the regex.
\end{pitfall}

\begin{pitfall}[document too large]
Symptom: saving a word fails with
\code{BsonMaximumSizeExceededException: ... exceeds 16793600 bytes}.
Cause: \code{images} holds base64 data URLs, and base64 inflates by a
third; five 3\,MB photos and the document is over the limit. Fix: store
images in MinIO under a \code{dictionary/\{userId\}/\{wordId\}/} prefix
and keep the keys in the document --- Chapter~35's rule, applied here.
Until then, cap the array in \code{WordDefinitionRequest} validation.
\end{pitfall}

\begin{pitfall}[the regex that never returns]
Symptom: one request pins a core at 100\,\% and the rest of the
service's queries queue behind it. Cause: a pattern like
\code{(a+)+\$} against a long string --- catastrophic backtracking.
\code{escapeRegex} prevents users from sending one; the risk returns
the day someone adds a ``power search'' that passes raw patterns, or
forgets the escape on a new endpoint. Defense in depth: a
\code{maxTimeMS} on the query (Spring Data's \code{@Meta(maxExecutionTimeMs)})
so a runaway pattern is killed server-side after a second.
\end{pitfall}

\section{What breaks at 3\,a.m.}

The MongoDB StatefulSet's node reboots; the pod is pending for four
minutes while the claim reattaches.

\begin{itemize}
  \item \textbf{dictionary-service} returns 500 on every endpoint. Its
  readiness probe (Chapter~17) does not include a database check, so
  the gateway keeps routing to it. Users see errors instead of a
  ``service unavailable'' page.
  \item \textbf{Everything else} is untouched. No other service holds a
  Mongo connection, no other service calls dictionary-service, and no
  Kafka topic depends on it. This is database-per-service working as
  designed: the blast radius of the outage is exactly one feature.
  \item \textbf{When the pod returns}, the driver's connection pool
  reconnects on its own; there is no restart needed and no cache to
  warm. Mongo's create-on-write means a \emph{wrong} claim would come up
  empty and silent (Chapter~8's pitfall) --- check \code{show dbs}
  before believing a ``data loss'' report.
\end{itemize}

The other 3\,a.m. is slower: the 5\,GB claim fills. Mongo does not fail
a single write cleanly at that point; it stalls, journals fail, and the
process may exit. Because documents can grow (images, \code{extra}),
there is no cap in the schema to prevent it. The alarm belongs on the
PVC's usage, not on the application.

\section{Enhancing the resolution}

Four changes, each complete, in the order they pay off.

\subsection{Page in the database}

Replace the list-returning query with a page-returning one. Spring Data
Mongo needs a separate count query for \code{@Query} pages:

\begin{lstlisting}[language=Java]
@Query(value = "{ 'createdByUserId': ?0, '$or': [ "
        + "{ 'word': { '$regex': ?1, '$options': 'i' } }, "
        + "{ 'tags': { '$regex': ?1, '$options': 'i' } } ] }",
       count = true)
long countByWordOrTag(String userId, String regex);

@Query("{ 'createdByUserId': ?0, '$or': [ "
        + "{ 'word': { '$regex': ?1, '$options': 'i' } }, "
        + "{ 'tags': { '$regex': ?1, '$options': 'i' } } ] }")
Page<WordDefinition> searchByWordOrTag(String userId, String regex,
                                       Pageable pageable);          (*@\co{1}@*)
\end{lstlisting}

\begin{lstlisting}[language=Java]
public Page<WordDefinitionResponse> searchWords(String userId,
        String query, Pageable pageable) {
    Pageable sorted = PageRequest.of(pageable.getPageNumber(),
            pageable.getPageSize(),
            Sort.by("word").ascending());                            (*@\co{2}@*)
    Page<WordDefinition> page = (query == null || query.isBlank())
            ? wordDefinitionRepository.findByCreatedByUserId(userId, sorted)
            : wordDefinitionRepository.searchByWordOrTag(userId,
                    escapeRegex(query.trim()), sorted);
    return toResponses(page, userId);                                (*@\co{3}@*)
}
\end{lstlisting}

\begin{description}
  \item[\co{1}] Spring Data appends \code{skip}, \code{limit} and
  \code{sort} to the query. The regex still scans the user's slice, but
  only twenty documents are fetched, serialized and mapped.
  \item[\co{2}] Sorting moves to the server. Case-insensitive order
  needs a collation on the query or the collection
  (\code{\{locale: "en", strength: 2\}}); add it to the index below so
  the sort can use it.
  \item[\co{3}] The link lookups become one batch per page, next.
\end{description}

\subsection{Batch the links}

\begin{lstlisting}[language=Java]
private Page<WordDefinitionResponse> toResponses(Page<WordDefinition> page,
                                                 String userId) {
    List<String> ids = page.map(WordDefinition::getId).toList();
    Map<String, List<String>> parents = new HashMap<>();
    Map<String, List<String>> children = new HashMap<>();
    for (WordLink l : wordLinkRepository.findByChildWordIdIn(ids)) {   (*@\co{1}@*)
        parents.computeIfAbsent(l.getChildWordId(), k -> new ArrayList<>())
               .add(l.getParentWordId());
    }
    for (WordLink l : wordLinkRepository
            .findByParentWordIdInOrderByPosition(ids)) {
        children.computeIfAbsent(l.getParentWordId(), k -> new ArrayList<>())
                .add(l.getChildWordId());
    }
    return page.map(w -> toResponse(w,
            parents.getOrDefault(w.getId(), List.of()),
            children.getOrDefault(w.getId(), List.of())));
}
\end{lstlisting}

\begin{description}
  \item[\co{1}] Two \code{\$in} queries per page instead of two per
  word: $2n + 1$ becomes $3$. Both use the existing indexes on
  \code{childWordId} and \code{parentWordId}.
\end{description}

\subsection{Index for the sort, and anchored search}

\begin{lstlisting}[language=Java]
// MongoIndexConfig
mongoTemplate.indexOps("word_definitions").ensureIndex(
        new CompoundIndexDefinition(
                new Document("createdByUserId", 1).append("word", 1))
        .collation(Collation.of("en").strength(2)));               (*@\co{1}@*)
\end{lstlisting}

\begin{description}
  \item[\co{1}] With this index, ``this user's words in alphabetical
  order'' is an index walk, so the unfiltered listing and every sorted
  page avoid an in-memory sort. And a \emph{prefix} search ---
  \code{\{word: \{\$regex: "\^{}clou"\}\}} with the same collation ---
  becomes a range seek on the same index. For a dictionary, prefix is
  usually what users mean when they type.
\end{description}

\subsection{A text index for word-anywhere search}

When substring-anywhere is genuinely required, stop using regex:

\begin{lstlisting}[language=Java]
// MongoIndexConfig: one text index per collection, weighted
mongoTemplate.indexOps("word_definitions").ensureIndex(
        new TextIndexDefinition.TextIndexDefinitionBuilder()
                .onField("word", 10F)                                (*@\co{1}@*)
                .onField("tags", 5F)
                .onField("meaning", 1F)
                .build());
\end{lstlisting}

\begin{lstlisting}[language=Java]
// DictionaryService
public Page<WordDefinitionResponse> textSearch(String userId,
        String terms, Pageable pageable) {
    TextCriteria text = TextCriteria.forDefaultLanguage()
            .matchingAny(terms.split("\\s+"));                       (*@\co{2}@*)
    Query q = TextQuery.queryText(text).sortByScore()
            .addCriteria(Criteria.where("createdByUserId").is(userId))
            .with(pageable);
    List<WordDefinition> hits = mongoTemplate.find(q, WordDefinition.class);
    long total = mongoTemplate.count(Query.of(q).skip(0).limit(0),
            WordDefinition.class);
    return toResponses(new PageImpl<>(hits, pageable, total), userId);
}
\end{lstlisting}

\begin{description}
  \item[\co{1}] A text index tokenizes and stems; \code{cloud} matches
  \code{clouds} and \code{clouded}, and a hit in \code{word} outranks
  the same hit in \code{meaning}. It does not match \emph{inside} a
  token (\code{lou} will not find \code{cloud}) --- that is the trade.
  \item[\co{2}] No regex, no escaping, no backtracking. The index is
  consulted first; \code{createdByUserId} filters the hits.
\end{description}

\begin{decision}[anchored prefix versus full text]
Prefix-anchored regex on the compound index answers ``words starting
with\ldots'' in microseconds and needs no new index type; it is the
right default for a type-ahead box. A text index answers ``entries about
\ldots'' with ranking and stemming, at the cost of a larger index and
no infix matching. Unanchored case-insensitive regex --- what the
project has --- answers ``anything containing\ldots'' and pays a scan
for it. The recommendation is prefix for the search box, text search
behind a separate ``search everywhere'' action, and regex retired.
Atlas Search (Lucene under Mongo's API) would give both at once, but
only on Atlas; on a self-hosted \code{mongod}, these two are the tools.
\end{decision}

\section{Outside the project}

Every document database has this chapter's shape: DynamoDB with its
partition-key-first access patterns and no substring search at all;
Couchbase with N1QL and a full-text index; Elasticsearch as the sidecar
that document stores grow when search becomes the product. Postgres's
\code{jsonb} plus \code{pg\_trgm} is the answer most teams should
consider before any of them --- it turns exactly this service's
\code{ILIKE '\%clou\%'} into a GIN index seek. The senior framing of
``SQL or NoSQL?'' is never the engine; it is the access patterns
written down first, then the engine whose indexes match them.

\section{Questions from real interviews}

\subsection*{Q: Dictionary search is slow. What do you check first?}

The plan, before anything else: \code{explain("executionStats")} on
the exact query the service sends, and two numbers from it ---
\code{totalDocsExamined} against \code{nReturned}, and which stage
carries the \code{filter}. In this service that shows an \code{IXSCAN}
on \code{createdByUserId} followed by a \code{FETCH} that examines
hundreds of documents per hit, because an unanchored case-insensitive
regex cannot use any index. Then look at what happens \emph{after} the
query: here, two link lookups per match and in-memory sorting and
paging. I would fix in that order --- page in the database, batch the
links with \code{\$in}, anchor the regex to a prefix on a compound
\code{\{createdByUserId, word\}} index with a case-insensitive collation
--- and re-read the plan after each step. If ``anywhere in the word''
is a real requirement, that is a text index, not a faster regex.

\subsection*{Q: When would you not pick MongoDB?}

When the access patterns are relational: data read in different
groupings than it is written, invariants that span entities
(a submission needs its assignment and its student to exist), or
reporting that joins across the domain. When the team already runs
Postgres well and the ``document'' is just a \code{jsonb} column away.
When multi-document transactions are routine rather than exceptional
--- Mongo has them, but on a replica set, with retry semantics you have
to code for. And when the data will be searched by substring at scale:
\code{pg\_trgm} is a one-line index; Mongo needs Atlas Search or an
external engine. The project's dictionary passes those tests because it
is per-user, nested, and joins nothing; the course domain fails all of
them, which is why it is in Postgres.

\subsection*{Q: Paginate a hundred thousand results without loading them.}

Never with \code{skip} at depth --- \code{skip(90000)} still walks
ninety thousand index entries. Use \emph{keyset} pagination: sort on an
indexed, unique-ish key (\code{word}, tie-broken by \code{\_id}) and
ask for \code{\{word: \{\$gt: lastWord\}\}} with \code{limit(20)}; each
page is an index seek regardless of depth. Return the last key as an
opaque cursor instead of a page number. Keep a count query only if the
UI truly needs a total, and cache it, because counting a filtered
regex is the same scan again. The project's \code{subList} paging is
the opposite of all this: it loads everything and returns twenty.

\subsection*{Q: One service uses Mongo, another Postgres. How do you keep their data consistent?}

You do not, transactionally --- there is no two-phase commit across
them and you should not want one. You keep them \emph{eventually}
consistent through events (Chapter~11): the owning service commits its
change and publishes a fact, the other consumes and updates its own
copy, and an outbox table in the owner makes ``committed'' and
``published'' the same thing. Where the project needs this today is the
user id: dictionary-service stores \code{createdByUserId} as a string
that auth-service issued, and never needs to join on it, so no
synchronization is required. If it ever needed the user's display name,
the answer is a \code{UserRegistered} consumer writing a tiny
\code{users} collection --- not a Feign call per word, and not a shared
database.

\subsection*{Q: The document shape changed. Migrate without downtime.}

Expand, migrate, contract. Deploy code that reads both shapes and
writes the new one (a \code{@Field} alias or a converter handles the
rename); run a background job that rewrites old documents in batches
with \code{updateMany} on a bounded \code{\_id} range, throttled,
idempotent, resumable; then, when a count of the old shape reaches
zero, deploy code that reads only the new shape. Mongo makes the first
step easy because nothing enforces the schema, and dangerous for the
same reason --- add a JSON Schema validator on the collection at the
end so the old shape cannot come back. For the \code{images} field in
this project the migration is also a data move: rewrite each data URL
into a MinIO object and replace the array element with its key.

\subsection*{Q: Model the word graph --- embed or reference, and what does traversal cost?}

Reference, as the project does, when edges have their own attributes
(\code{position}) and both directions are queried; embed a
\code{childIds} array when the graph is shallow, mostly read, and
parent-to-child is the only direction that matters. Traversal cost is
the honest part: with references and naive recursion it is two round
trips per node --- 4{,}000 for a 2{,}000-word graph --- which is what
\code{buildGraphNode} does. The fixes are, in order, loading all edges
for the user in one query and walking in memory; or a single
\code{\$graphLookup} with \code{maxDepth}; or embedding. Cycles need a
visited set either way; the project has one. And say the trade-off out
loud: embedding turns one link creation into two document writes with
no transaction on a standalone node.

\begin{quiz}
\begin{enumerate}
  \item From the \code{explain} output, \code{totalDocsExamined} is 412
  and \code{nReturned} is 3. Name the stage doing the rejecting, say
  why the \code{createdByUserId} index cannot help it, and state the
  one change to the pattern that would let an index on \code{word}
  serve it.
  \item \code{searchWords} with \code{page=50\&size=20} on a user with
  3{,}000 words: how many documents does Mongo return, how many
  \code{word\_links} queries run, and where is the page built?
  \item \code{deleteWord} issues two writes with no transaction.
  Explain why a transaction is unavailable on this deployment, what a
  crash between the writes leaves behind, and which line of
  \code{buildGraphNode} tolerates it.
  \item A user attaches six photos to a word and the save fails. Give
  the exception, the arithmetic that produced it, and the storage rule
  from Chapter~35 that fixes it.
  \item With a text index on \code{word}, \code{tags} and
  \code{meaning}, the query ``lou'' returns nothing although
  ``cloud'' exists. Explain, and say which index type answers it.
  \item MongoDB is down for four minutes at 3\,a.m. List what fails,
  what does not, and the one probe change that would turn 500s into a
  clean ``not ready''.
\end{enumerate}
\end{quiz}
